%Schriftart
\usepackage[utf8]{inputenc}
\usepackage[T1]{fontenc}
\usepackage[ngerman]{babel}
\usepackage{lmodern}

%schriftgröße der Überschriften
\setkomafont{chapter}{\Large}
\setkomafont{section}{\large}
\setkomafont{subsection}{\normalsize}

%zeilenabstand der itemizes
\usepackage{enumitem}
\setlist[itemize]{itemsep=0pt, parsep=0pt, topsep=0pt, partopsep=0pt}

%beginn fußnoten mit leerzeichen
\deffootnote{1em}{1em}{\makebox[1em][r]{\thefootnotemark\ }}

%normale Fußnoten mit Punkt schließen lassen, wenn noch kein '.', '!', '?' vorhanden
\let\orgfootnote\footnote
\newcommand\myautodot{%
  \ifthenelse{\the\spacefactor>\sfcode`,}{}{.}%
}
\renewcommand\footnote[2][\empty]{%
  \ifx#1\empty%
    \orgfootnote{#2\myautodot}%
  \else%
    \orgfootnote[#1]{#2\myautodot}%
  \fi%
}

%anpassung seitenränder
\usepackage[
    left=25mm,   % linker Rand
    right=25mm,  % rechter Rand
    top=25mm,    % oberer Rand
    bottom=20mm  % unterer Rand
]{geometry}

%Festlegen des Seitenlayouts
\usepackage[onehalfspacing]{setspace} %Zeilenabstand
%alternativ \linespread{1.5}
%BCOR ist der Binderand, DIV definiert den Satzspiegel (siehe KOMA-script Dokumentation S. 30)
\KOMAoptions{BCOR=0.5cm, DIV=11, parskip=half}
%Absätze werden über parskip geregelt. Die Werte sind parskip=true/false/half

%Abstand vor Kapitelbeginn
\RedeclareSectionCommand[
  beforeskip=0pt,
  afterindent=false
]{chapter}

%Schriftart der Überschriften serif in bold
\addtokomafont{disposition}{\rmfamily}
%Schriftart der Überschriften in serif ohne bold
%\setkomafont{disposition}{\normalfont}

%Für das Einfügen eines Glossars: siehe https://en.wikibooks.org/wiki/LaTeX/Glossary

%stützt die Schnittstelle von float und listings zu tocbasic (genutzt von KOMA-script)
\usepackage{scrhack}

%Zum Auskommentieren von größeren Blöcken
\usepackage{comment}
%Als Platzhalter
\usepackage{lipsum}
%Einfachere Verwendung von Anführungszeichen
\newcommand{\qt}[1]{\glqq#1\grqq{}}
%Angabe von Bildquellen
\newcommand{\quelle}[1]{\\\vspace{.1cm}{Quelle: #1}}

%Bilder
\usepackage{graphicx}
\usepackage{float}
\graphicspath{{./abbildungen/}}
\usepackage[export]{adjustbox}
\usepackage{array}


%Einfügen von Unformatiertem typewriter Text
\usepackage{verbatim}
%über mehrere Zeilen mit \begin{Verbatim}[breaklines=true]\end{Verbatim}
\usepackage{fancyvrb}


%Alles was das Formelherz begehrt
\usepackage{amstext} 
\usepackage{amsmath}
\usepackage{amssymb}

%Für die Literaturverzeichnisse
\usepackage{csquotes}
%\usepackage[backend=biber, bibstyle=authoryear, citestyle=apa]{biblatex}
%Für Hauptname der Quelle in anderer Schriftart nehme man 'iso-authoryear'. (nutzt dreprecated Feld)




%Literaturverzeichnis exakt nach ISO 690-3
%Eine Alternative Erstellung von Literaturverzeichnissen über natbib (letzte Version von 2010)
%WICHTIG: bei einem Wechsel müssen alle Temporären Dateien von LaTex gelöscht werden!
%\usepackage{natbib}
%\usepackage{multibib} %Unterteilung in Internet- und Printmedien
%\newcites{online}{Internetquellenverzeichnis} 
%Statt \cite{} muss für Internetquellen nun \citeonline{} verwendet werden!



%Anklickbare Links und \url{}
\usepackage{hyperref}
%Zeilenumbrüche in URLs, standard url Paket wird von hyperref geladen, Befehl \url{} bleibt gleich
\usepackage{microtype}[final]
\urlstyle{same}

%glossar
\usepackage[acronym,hyperfirst=true,toc=true]{glossaries}
%\hypersetup{pageanchor=true}
\makeglossaries
%https://de.overleaf.com/learn/latex/Glossaries infos über Erstellung glossareintraäge und aufrufe dieser
%WICHTIG! wenn Änderungen an den Glossareinträgen vorgenommen werden,im Terminal hintereinander das hier iengeben:
%erstmal .glo datei lsöchen, dann:
%pdflatex Seminarfacharbeit.tex
%makeglossaries Seminarfacharbeit
%pdflatex Seminarfacharbeit.tex
%nur dann funktioniert das Glossar auch in der PDF-Datei
\begin{comment}
    zum kopieren:
    \newglossaryentry{}
    {}
\end{comment}

\newglossaryentry{silber-silberchlorid-elektrode}
    {
    name=Silber-Silberchlorid-Elektrode,
    text={\textit{Silber-Silberchlorid-Elektrode}},
    first={\textit{Silber-Silberchlorid-Elektrode}},
    description={Die Silber-Silberchlorid-Elektrode ist eine häufig verwendete Elektrode in der Elektroenzephalographie und besteht aus einem Silberdraht, der mit einer Schicht aus Silberchlorid überzogen ist}
    }

\newglossaryentry{berger-effekt}
{
    name=Berger-Effekt,
    text={\textit{Berger-Effekt}},
    first={\textit{Berger-Effekt}},
    description={Der Berger-Effekt bezeichnet eine Blockade der Alpha-Wellen im EEG und ist nachweisbar, in dem ein Proband, der wach und entspannt ist, die Augen öffnet. Die normalerweise bei geschlossenen Augen sichtbaren Alpha-Wellen verschwinden und werden durch Beta-Wellen ersetzt. Beim Schließen der Augen kehrt der Alpha-Rhythmus zurück}
}

\newglossaryentry{brain-computer-interfaces}
{
    name=Brain-Computer-Interfaces,
    text={\textit{Brain-Computer-Interfaces}},
    first={\textit{Brain-Computer-Interfaces}},
    description={Brain-Computer-Interfaces (BCI) sind Systeme, die es ermöglichen, Informationen direkt aus dem Gehirn zu extrahieren und in Befehle für externe Geräte umzuwandeln}
}

\newglossaryentry{elektrolyt}
{
    name=Elektrolyt,
    text={\textit{Elektrolyt}},
    first={\textit{Elektrolyt}},
    description={Ein Elektrolyt ist eine Substanz, die in der Lage ist, elektrische Ladungen zu leiten, wenn sie als Lösung oder im geschmolzenem Zustand vorliegt}
}

\newglossaryentry{analog-digital-converter}
    {name=Analog-Digital-Converter,
    text={\textit{Analog-Digital-Converter}},
    first={\textit{Analog-Digital-Converter}},
    description={Ein Analog-Digital-Converter, auf deutsch Analog-Digital-Wandler, ist ein elektronisches Bauteil, das analoge Signale in einen digitalen Datenstrom, der dann weiterverarbeitet oder gespeichert werden kann}
}

\newglossaryentry{notch}{
    name=Notch-Filter,

    description={Filter, der gezielt eine bestimmte Frequenz (z.B. 50 Hz Netzrauschen) unterdrückt}
}

\newglossaryentry{ica}{
    name=ICA,
    description={Independent Component Analysis, mathematisches Verfahren zur Trennung überlagerter Signalquellen im EEG}
}

\newglossaryentry{bandpass}{
    name=Bandpassfilter,
    description={Filter, der nur Signale innerhalb eines bestimmten Frequenzbereichs durchlässt}
}

\newglossaryentry{cutoff}{
    name=Cut-off-Frequenz,
    description={Grenzfrequenz, ab der ein Filter Signale durchlässt oder unterdrückt}
}

\newglossaryentry{p300}{
    name=P300,
    description={ERP-Komponente, die etwa 300 ms nach einem seltenen oder bedeutsamen Stimulus auftritt}
}

\newglossaryentry{eog}{
    name=EOG,
    description={Elektrookulogramm, Messung von Augenbewegungen zur Identifikation von Artefakten im EEG}
}

\newglossaryentry{eventmarker}{
    name=Ereignismarker,
    description={Zeitliche Markierung im EEG-Signal zur Kennzeichnung eines Stimulus oder Ereignisses}
}

\newglossaryentry{oddball}{
    name=Oddball-Aufgabe,
    description={Experimentelles Paradigma, bei dem seltene Zielreize zwischen häufigen Standardreizen präsentiert werden, um ERPs zu untersuchen}
}




%Formatierung und Import Bibliography

\usepackage[affil-it]{authblk}	
\usepackage[datamodel=datamodel,backend=biber,style=authoryear,autocite=footnote,maxnames=20,minnames=10]{biblatex}
\addbibresource{literatur.bib}

% slash zwischen namen 
\DeclareDelimFormat{multinamedelim}{\addslash\addspace}
\DeclareDelimFormat{finalnamedelim}{\addslash\addspace}
\DeclareDelimFormat{multinamedelim}{\addspace\addslash\addspace}
\DeclareDelimFormat{finalnamedelim}{\addspace\addslash\addspace}
%Vor und nach name richtige reinfolge
\DeclareNameAlias{sortname}{family-given}   
\DeclareNameAlias{default}{family-given}    

\DefineBibliographyStrings{ngerman}{
  urlseen = {Abgerufen am }
}

%\DeclareFieldFormat{note}{\mkbibbrackets{\bibstring{urlseen}\addcolon\space#1}}

\DeclareFieldFormat{urldate}{\mkbibbrackets{\bibstring{urlseen}\addcolon\space#1}}

\DeclareFieldFormat{lasteditdate}{\mkbibparens{#1}}

\DeclareFieldFormat{year}{\mkbibparens{#1}}

\DeclareFieldFormat{volume}{#1}






%citation commands
\DeclareCiteCommand{\vgl}[\mkbibfootnote]{%
	\usebibmacro{citeindex}%
	\iffieldundef{prenote}
		{\printtext{Vgl\adddot\space}}
		{\printfield{prenote}\setunit{\prenotedelim}}%
}
{%
	% Ausgabe abhängig vom Eintragstyp
	\ifentrytype{onlinedokument}
		{\printnames{labelname}\addspace\bibhyperref{\printurldate}}%
		{%
					\ifentrytype{onlinelexikon}
						{\printfield{journaltitle}\adddot\addspace\bibhyperref{\printfield{title}}}%
						{%
							\ifentrytype{zeitschriftartikel}
								{\printnames[family]{labelname}\addspace\bibhyperref{\printfield{year}}}%
								{\printnames[family]{labelname}\addspace\bibhyperref{\printfield{year}}}%
						}%
		}%
}
{\multicitedelim}
{%
	\iffieldundef{postnote}
		{\adddot}
		{\addcolon\addspace\printfield{postnote}\adddot}%
	\nopunct%
}

\DeclareCiteCommand{\vglHier}[]{%
	\usebibmacro{citeindex}%
	\iffieldundef{prenote}
		{\printtext{Vgl\adddot\space}}
		{\printfield{prenote}\setunit{\prenotedelim}}%
}
{%
	% Ausgabe abhängig vom Eintragstyp
	\ifentrytype{onlinedokument}
		{\printnames{labelname}\addspace\bibhyperref{\printurldate}}%
		{%
					\ifentrytype{onlinelexikon}
						{\printfield{journaltitle}\adddot\addspace\bibhyperref{\printfield{title}}}%
						{%
							\ifentrytype{zeitschriftartikel}
								{\printnames[family]{labelname}\addspace\bibhyperref{\printfield{year}}}%
								{\printnames[family]{labelname}\addspace\bibhyperref{\printfield{year}}}%
						}%
		}%
}
{\multicitedelim}
{%
	\iffieldundef{postnote}
		{\adddot}
		{\addcolon\addspace\printfield{postnote}\adddot}%
	\nopunct%
}
	


\DeclareCiteCommand{\ebd}[\mkbibfootnote]
{\usebibmacro{citeindex}%
\iffieldundef{prenote}
{\printtext{Vgl\adddot}}
{\printfield{prenote}\setunit{\prenotedelim}}%
}
{
	\printtext{ebd\adddot}
	\iffieldundef{postnote}
	{}
	{\addspace\printfield{postnote}\adddot}
}
{}
{}
















%Einbinden von Quelltext in Python (https://tex.stackexchange.com/questions/83882/how-to-highlight-python-syntax-in-latex-listings-lstinputlistings-command)
% Default fixed font does not support bold face
\DeclareFixedFont{\ttb}{T1}{txtt}{bx}{n}{12} % for bold
\DeclareFixedFont{\ttm}{T1}{txtt}{m}{n}{12}  % for normal

% Custom colors
\usepackage{color}
\definecolor{deepblue}{rgb}{0,0,0.5}
\definecolor{deepred}{rgb}{0.6,0,0}
\definecolor{deepgreen}{rgb}{0,0.5,0}

\usepackage{listings}

% Python style for highlighting
\newcommand\pythonstyle{\lstset{
language=Python,
basicstyle=\ttm,
otherkeywords={self},             % Add keywords here
keywordstyle=\ttb\color{deepblue},
emph={MyClass,__init__},          % Custom highlighting
emphstyle=\ttb\color{deepred},    % Custom highlighting style
stringstyle=\color{deepgreen},
frame=tb,                         % Any extra options here
showstringspaces=false            % 
}}


% Python environment
\lstnewenvironment{python}[1][]
{
\pythonstyle
\lstset{#1}
}
{}

% Python for external files
\newcommand\pythonexternal[2][]{{
\pythonstyle
\lstinputlisting[#1]{#2}}}

% Python for inline
\newcommand\pythoninline[1]{{\pythonstyle\lstinline!#1!}}

\usepackage{eurosym}
